% ==================================================================
% id:       lim-043
% topic:    גבולות ורציפות של פונקציות
% subtopic: אינסוף כגבול וגבול באינסוף
% solution: full
% ==================================================================
\textbf{
בכל אחד מהסעיפים הבאים, השתמשו במשפט היינה וקבעו האם קיים הגבול. \\
נמקו את תשובתכם/ן:}

\begin{enumerate}\bMs{2}

\item $\ds{\lim_{x \to 0} e^{\frac{1}{x}}}$
\ifsol
\hint{בדקו לאן שואפת $f(a_n)$ עבור סדרות $a_n\to0$ מימין ומשמאל, והשוו.}
\begin{solbox}
נראה כי הגבול לא קיים.
תהיינה
$(a_n)_{n=1}^\infty$
ו-
$(b_n)_{n=1}^\infty$
שתי סדרות כך ש-
$$,\lim_{n\to\infty}a_n=\lim_{n\to\infty}b_n=0$$
ו-
$a_n>0$, $b_n<0$
לכל
$n\in\mathbb{N}$.
מתקיים
$$,\displaystyle\lim_{n\to\infty}f(a_n)=\lim_{n\to\infty}e^{\frac{1}{a_n}}
=\infty$$
מאחר והגבול במעריך הוא אינסוף.
מצד שני
$$,\displaystyle\lim_{n\to\infty}f(b_n)=\lim_{n\to\infty}e^{\frac{1}{b_n}}=0$$
מאחר והגבול במעריך הוא מינוס אינסוף.
התקבלו שתי תוצאות שונות, משמעו אין גבול.
לדוגמה, אפשר לבחור $a_n=\frac1n$ ו-$b_n=-\frac1n$. אז
$$f(a_n)=e^{n}\xrightarrow[n\to\infty]{}\infty,\qquad f(b_n)=e^{-n}\xrightarrow[n\to\infty]{}0$$
\end{solbox}
\fi

\item $\ds{\lim_{x \to 0} \arctan\lr{ \frac{1}{x} } }$
\ifsol
\hint{מה קורה ל-$\frac1x$ כאשר מתקרבים ל-$0$ מכל צד, ומה ידוע על $\arctan$ בפלוס ומינוס אינסוף?}
\begin{solbox}
נזכיר איך נראה הגרף של $\arctan$: היא מוגדרת לכל $x\in\mathbb{R}$, עולה ממש, והתמונה שלה היא $\lr{-\frac\pi2,\frac\pi2}$. כאשר $x\to\infty$ הגרף מתקרב לישר $y=\frac\pi2$, וכאשר $x\to-\infty$ לישר $y=-\frac\pi2$. כלומר
$$\lim_{x\to\infty}\arctan(x)=\frac\pi2,\qquad \lim_{x\to-\infty}\arctan(x)=-\frac\pi2$$
\qfig{lim-043b}
שוב נראה כי הגבול לא קיים. נבחר שתי סדרות - הראשונה שמתכנסת ל-
$0$
מלמעלה, והשניה - מלמטה.
לדוגמא,
$a_n=-\tfrac{1}{n}$
ו-
$b_n=\tfrac{1}{n}$.
מתקיים
\begin{align*}
\displaystyle\lim_{n\to\infty}f(a_n)
&=\ \displaystyle\lim_{n\to\infty}\arctan(-n)\\
&=\ -\frac{\pi}{2},
\end{align*}
מצד שני
\begin{align*}
\displaystyle\lim_{n\to\infty}f(b_n)
&=\ \displaystyle\lim_{n\to\infty}\arctan(n)\\
&=\ \frac{\pi}{2},
\end{align*}
התקבלו שני גבולות שונים, לכן הגבול המבוקש לא קיים
(ע"פ היינה).
\end{solbox}
\fi

\item $\ds{\lim_{x\to0}\frac{3+2e^{\frac1x}}{1+e^{\frac1x}}}$
\ifsol
\hint{בחרו סדרה ששואפת ל-$0$ מימין וסדרה ששואפת ל-$0$ משמאל. מה קורה ל-$e^{\frac1x}$ בכל אחת מהן?}
\begin{solbox}
נראה כי הגבול לא קיים. נבחר $a_n=\frac1n$ ו-$b_n=-\frac1n$. שתי הסדרות שואפות ל-$0$, ואיבריהן שונים מ-$0$.
\begin{itemize}
\item כאן $f(a_n)=\frac{3+2e^{n}}{1+e^{n}}$. נחלק את המונה ואת המכנה ב-$e^n>0$:
\[f(a_n)=\frac{3e^{-n}+2}{e^{-n}+1}\]
מתקיים $\ds\lim_{n\to\infty}e^{-n}=0$, ולכן לפי אריתמטיקה של גבולות
\[\lim_{n\to\infty}f(a_n)=\frac{0+2}{0+1}=2\]
\item כאן $f(b_n)=\frac{3+2e^{-n}}{1+e^{-n}}$, ולפי אריתמטיקה של גבולות
\[\lim_{n\to\infty}f(b_n)=\frac{3+0}{1+0}=3\]
\end{itemize}
התקבלו שני גבולות שונים, ולכן לפי משפט היינה הגבול $\ds\lim_{x\to0}\frac{3+2e^{\frac1x}}{1+e^{\frac1x}}$ אינו קיים.
\end{solbox}
\fi
\eMs
\end{enumerate}
